% 图契约：由能量证明互感上限；图中的端口、几何、箭头和方程对应本题。
% 版面：沿用前三模块的几何图+推导分区；目标印宽17cm，简图至中等密度。
% 布局：左侧真实结构，右侧关键关系；真实电路边不作为架构流程箭头处理。
\documentclass[tikz,border=7pt]{standalone}
\usepackage{fontspec,xeCJK,amsmath,amssymb}
\IfFontExistsTF{Times New Roman}{\setmainfont{Times New Roman}}{\setmainfont{TeX Gyre Termes}}
\IfFontExistsTF{Heiti SC}{\setCJKmainfont[AutoFakeBold=1.5]{Heiti SC}}{\setCJKmainfont[AutoFakeBold=1.5]{FandolHei-Regular}}
\usetikzlibrary{arrows.meta,bending,calc,patterns,positioning,decorations.pathmorphing}
\definecolor{ink}{HTML}{243745}
\definecolor{blue}{HTML}{275DAB}
\definecolor{green}{HTML}{177C63}
\definecolor{red}{HTML}{BD4935}
\definecolor{orange}{HTML}{AA6C16}
\definecolor{purple}{HTML}{715393}
\tikzset{every node/.style={font=\fontsize{9.3}{12.5}\selectfont,text=ink},
arrow/.style={-{Stealth[length=5pt 1.5,width'=0pt 3]},shorten >=2pt,shorten <=1pt,line width=1pt},
arrow short/.style={-{Stealth[length=3pt 1.0,width'=0pt 2.5]},shorten >=1pt,shorten <=0pt,line width=.8pt},
axis/.style={arrow,draw=ink!60,line width=.6pt},
vec/.style={arrow,draw=blue},
wire/.style={draw=ink!75,line width=.8pt},
guide/.style={dashed,draw=ink!45,line width=.5pt},
card/.style={rounded corners=2pt,draw=blue!65,fill=blue!4,align=center,text width=6.5cm,inner xsep=9pt,inner ysep=7pt},
title/.style={font=\fontsize{10.5}{14}\selectfont\bfseries,text width=15cm,align=center},
label/.style={font=\fontsize{8.4}{11}\selectfont,align=center},
note/.style={label,text width=6.1cm}}
\begin{document}
\begin{tikzpicture}[x=1cm,y=1cm]
\path[use as bounding box] (-8,-3.5) rectangle(8,3.65);
\node[title] at (0,3.25) {M5-094\quad 由能量证明互感上限};
\begin{scope}[shift={(-4.05,0)}]
\draw[wire](-1.5,1.8)--(-1.5,1.1);\draw[wire](-1.5,-1.1)--(-1.5,-1.8);\fill[orange](-1.77,1.1)circle(2.5pt);\node[above]at(-1.5,2.05){$+\ v_1$};\node[below]at(-1.5,-2.05){$-$};\draw[wire](-1.5,1.1)..controls(-0.85,1.1)and(-0.85,0.55)..(-1.5,0.55);\draw[wire](-1.5,0.55)..controls(-0.85,0.55)and(-0.85,0.0)..(-1.5,0.0);\draw[wire](-1.5,0.0)..controls(-0.85,0.0)and(-0.85,-0.55)..(-1.5,-0.55);\draw[wire](-1.5,-0.55)..controls(-0.85,-0.55)and(-0.85,-1.1)..(-1.5,-1.1);\draw[arrow short,red](-2.2,1.8)--(-2.2,.7)node[midway,left]{$i_1$};\draw[wire](1.5,1.8)--(1.5,1.1);\draw[wire](1.5,-1.1)--(1.5,-1.8);\fill[orange](1.23,1.1)circle(2.5pt);\node[above]at(1.5,2.05){$+\ v_2$};\node[below]at(1.5,-2.05){$-$};\draw[wire](1.5,1.1)..controls(2.15,1.1)and(2.15,0.55)..(1.5,0.55);\draw[wire](1.5,0.55)..controls(2.15,0.55)and(2.15,0.0)..(1.5,0.0);\draw[wire](1.5,0.0)..controls(2.15,0.0)and(2.15,-0.55)..(1.5,-0.55);\draw[wire](1.5,-0.55)..controls(2.15,-0.55)and(2.15,-1.1)..(1.5,-1.1);\draw[arrow short,red](0.8,1.8)--(0.8,.7)node[midway,left]{$i_2$};\node at(0,.2){$M$};\end{scope}
\node[card,draw=blue!65,fill=blue!4] at (3.9,1.8) {$\begin{aligned}W_m={}&\frac12L_1(i_1+Mi_2/L_1)^2\\&+\frac12(L_2-M^2/L_1)i_2^2\end{aligned}$};
\node[card,draw=green!65,fill=green!4] at (3.9,0) {$L_2-M^2/L_1\ge0$};
\node[card,draw=orange!65,fill=orange!4] at (3.9,-1.8) {$M^2\le L_1L_2$};
\end{tikzpicture}\end{document}